\begin{figure}[t]
\centering
\begin{tikzpicture}[
  font=\small,
  comp/.style={draw, rounded corners=2pt, inner sep=5pt, align=left, text width=33mm, font=\footnotesize},
]

% Sjølve likninga
\node (eq) at (0,0) {$\displaystyle \mathrm{Form} \;=\; SG\!\left(\nabla_{C(A)}\, L(c,t)\right) \,\cap\, \lozenge$};

% Tre forklaringsboksar under
\node[comp] (sg) at (-40mm,-22mm) {$SG$ = formgrammatikk};
\node[comp] (l)  at (0,-22mm)     {$L$ = tilpassingslandskap (samtidige trykk)};
\node[comp] (lo) at (40mm,-22mm)  {$\lozenge$ = kognitiv lyskjegle (agentens horisont)};

% Pilar frå boks til komponent i likninga
\draw[-{Latex[length=2mm]}] (sg.north) .. controls (-40mm,-8mm) and (-30mm,-2mm) .. (eq.south west);
\draw[-{Latex[length=2mm]}] (l.north) -- (eq.south);
\draw[-{Latex[length=2mm]}] (lo.north) .. controls (40mm,-8mm) and (30mm,-2mm) .. (eq.south east);

\end{tikzpicture}
\caption{\textsc{Den generative likninga, lagdelt.}
Forma kjem ut av tre lag. Formgrammatikken $SG$ avgrensar kva som i det heile
kan setjast saman. Gradienten $\nabla_{C(A)} L(c,t)$ peikar mot den retninga
der tilpassingslandskapet $L$ --- summen av samtidige trykk på tida $t$ ---
fell raskast. Snittet med lyskjegla $\lozenge$ avgrensar resultatet til det som
ligg innanfor agentens kognitive horisont: berre former agenten kan tenkje seg,
blir til.}
\label{fig:formlaere-likning}
\end{figure}
